% ==============================================================================
% פרק 2: סקירת ספרות ומסגרת תיאורטית: פסיכולוגיה חברתית ומרחבים דיגיטליים
% ==============================================================================
\chapter{סקירת ספרות ומסגרת תיאורטית: פסיכולוגיה חברתית במרחב הדיגיטלי}
\label{chap:literature_review}

\section{מבוא לפרק}
פרק זה מבסס את התשתית המדעית והרעיונית של פלטפורמת \projectname{}. כדי להבטיח שהארכיטקטורה מרובת הסוכנים תישען על עקרונות פסיכולוגיים מבוססים ולא על השערות אינטואיטיביות, הפרק מנתח לעומק את תיאוריות הליבה של פסיכולוגיה חברתית ומדיה דיגיטלית:
\begin{itemize}
    \item \textbf{תיאוריית ההשוואה החברתית של פסטינגר}~\cite{festinger1954theory}, מודל הפרוקסי של סאלס וחוב'~\cite{suls2002social}, מנגנוני הניגוד וההטמעה, והאופן שבו רשתות אלגוריתמיות מפרות את האיזון הטבעי של תהליכים אלו;
    \item \textbf{תיאוריית הפער העצמי של היגינס}~\cite{higgins1987self} ותיאוריית ההחפצה העצמית במרחב הדיגיטלי~\cite{perloff2014social};
    \item \textbf{סקירה שיטתית של ספרות אמפירית} העוסקת בחשיפה פסיבית לרשתות, שיווק משפיענים ורווחה סובייקטיבית;
    \item \textbf{מודל מושגי רב-משתני} הממפה את הגורמים המתווכים והממתנים בין חשיפה דיגיטלית לרווחה נפשית;
    \item \textbf{כלי מדידה פסיכומטריים מתוקפים} (כגון \textenglish{RSES} ו-\textenglish{PANAS}) והמרתם למשתני מצב וקטוריים במערכות בינה מלאכותית;
    \item \textbf{סוכנים גנרטיביים כמודל לסימולציה אנושית} בהשראת מאמר הדגל של פארק ועמיתיו~\cite{park2023generative}.
\end{itemize}

\section{תיאוריית ההשוואה החברתית של פסטינגר במרחב הדיגיטלי}
\label{sec:lit_sct}
בשנת 1954 ניסח לאון פסטינגר (\textenglish{Leon Festinger}) את תיאוריית ההשוואה החברתית (\textenglish{Social Comparison Theory - SCT})~\cite{festinger1954theory}, שהפכה לאחת מאבני היסוד המרכזיות של הפסיכולוגיה החברתית. הנחת המוצא של פסטינגר היא כי בקרב בני אדם קיים דחף פנימי מולד (\textenglish{Innate Drive}) להעריך את יכולותיהם, מעמדם והישגיהם. כאשר לא קיימים מדדים אובייקטיביים, הפרט משווה את עצמו לאחרים בסביבתו בשני כיוונים עיקריים:
\begin{enumerate}
    \item \textbf{השוואה כלפי מטה (\textenglish{Downward Social Comparison})}: השוואה עצמית ליחידים הנתפסים כנמצאים במצב נחות. תהליך זה מונע מהאדרה עצמית (\textenglish{Self-Enhancement}) ומשמש לוויסות רגשי חיובי ולהגנת הערך העצמי.
    \item \textbf{השוואה כלפי מעלה (\textenglish{Upward Social Comparison})}: השוואה לעולים על הפרט בממדי יופי, עושר, אושר או הצלחה. מודל הפרוקסי של סאלס ועמיתיו~\cite{suls2002social} מדגים כי השוואה זו מתפצלת בין \textbf{אפקט הטמעה והשראה} (כאשר היעד נתפס כבר-השגה) לבין \textbf{אפקט ניגוד ושחיקה} (כאשר היעד נתפס כבלתי-מושג, מה שמעורר קנאה, תסכול ודכדוך).
\end{enumerate}

בעוד שבסביבה הפיזית תהליכי ההשוואה מאוזנים על ידי מפגשים יומיומיים מגוונים, סביבות הרשת המודרניות מפעילות \textbf{שיבוש אלגוריתמי עמוק}: אסימטריית אצירה (\textenglish{Highlight Reels}), ריטוש תמונות ומסננים דיגיטליים, ואלגוריתמי אופטימיזציית מעורבות המקדמים תכנים מקצינים. כתוצאה מכך, הפיד של צעירים מוטה באופן מובנה להשוואה כלפי מעלה ולאפקט הניגוד~\cite{appel2015social, appel2016interplay}.

\begin{figure}[H]
\centering
\begin{LTR}
\begin{tikzpicture}[
    scale=0.62, every node/.style={transform shape},
    corenode/.style={rectangle, draw=DarkPurple, fill=DarkPurple, text=white, rounded corners=3mm, line width=1.3pt, align=center, inner sep=6pt, text width=3.4cm, font=\small\bfseries\sffamily},
    dirnode/.style={rectangle, draw=RoyalPurple, fill=LavenderCard, rounded corners=2.5mm, line width=1.1pt, align=center, inner sep=5pt, text width=3.2cm, font=\footnotesize\sffamily},
    effectnode/.style={rectangle, draw=PurpleAccent, fill=white, rounded corners=2.5mm, line width=1.0pt, align=center, inner sep=5pt, text width=3.2cm, font=\footnotesize\sffamily},
    outcomepos/.style={rectangle, draw=SlateMuted, fill=white, rounded corners=2.5mm, line width=1.0pt, align=center, inner sep=5pt, text width=3.2cm, font=\footnotesize\sffamily},
    outcomeneg/.style={rectangle, draw=AlertRose, fill=AlertRose!15, rounded corners=2.5mm, line width=1.2pt, align=center, inner sep=5pt, text width=3.2cm, font=\footnotesize\sffamily},
    arrow/.style={-{Latex[length=2.5mm]}, line width=1.0pt, draw=SlateGray},
    highlightarrow/.style={-{Latex[length=2.8mm]}, line width=1.4pt, draw=AlertRose}
]

\node[corenode] (core) {\textbf{Leon Festinger (1954)}\\\scriptsize Social Comparison Theory\\\scriptsize Innate Drive for Evaluation};
\node[dirnode, above right=0.6cm and 1.3cm of core] (upward) {\textbf{Upward Comparison}\\\scriptsize Superior Targets\\\scriptsize (Lifestyle, Beauty, Wealth)};
\node[dirnode, below right=0.6cm and 1.3cm of core] (downward) {\textbf{Downward Comparison}\\\scriptsize Inferior Targets\\\scriptsize (Self-Enhancement)};
\node[outcomepos, right=1.3cm of downward] (down_outcome) {\textbf{Reassurance \& Defense}\\\scriptsize \textenglish{$\uparrow S_t$ (Self-Esteem)}\\\scriptsize Anxiety Reduction};
\node[effectnode, above right=0.2cm and 1.3cm of upward] (contrast) {\textbf{Contrast Effect}\\\scriptsize Unattainable Target\\\scriptsize \textbf{Algorithmic Default!}};
\node[effectnode, below right=0.2cm and 1.3cm of upward] (assim) {\textbf{Assimilation Effect}\\\scriptsize Attainable Target\\\scriptsize (Inspiration \& Growth)};
\node[outcomeneg, right=1.2cm of contrast] (contrast_out) {\textbf{Affective Depletion}\\\scriptsize Envy, Resentment\\\scriptsize \textenglish{$\uparrow$ Self-Discrepancy}\\\scriptsize \textenglish{$\downarrow$ State Self-Esteem}};
\node[outcomepos, right=1.2cm of assim] (assim_out) {\textbf{Constructive Motivation}\\\scriptsize Optimism, Growth};

\draw[highlightarrow] (core) -- node[above, sloped, font=\scriptsize\sffamily, text=AlertRose] {Algorithmic Bias} (upward);
\draw[arrow] (core) -- (downward);
\draw[arrow] (downward) -- (down_outcome);
\draw[highlightarrow] (upward) -- node[above, sloped, font=\scriptsize\sffamily, text=AlertRose] {Hyper-Idealization} (contrast);
\draw[arrow] (upward) -- (assim);
\draw[highlightarrow] (contrast) -- (contrast_out);
\draw[arrow] (assim) -- (assim_out);
\end{tikzpicture}
\end{LTR}
\caption{מודל תיאוריית ההשוואה החברתית של פסטינגר: כיווניות, מנגנוני ניגוד והטמעה, וההטיה האלגוריתמית}
\label{fig:sct_model}
\end{figure}

\section{תיאוריית הפער העצמי והחפצה עצמית במרחב הדיגיטלי}
\label{sec:lit_self_discrepancy}
כדי להבין את המנגנון הקוגניטיבי המתווך בין חשיפה להשוואה כלפי מעלה לבין הופעת מצוקה רגשית, מודל \projectname{} משלב את תיאוריית הפער העצמי (\textenglish{Self-Discrepancy Theory}) של אדוארד טורי היגינס (\textenglish{E. Tory Higgins})~\cite{higgins1987self}. היגינס מגדיר שלושה מרחבי עצמי:
\begin{enumerate}
    \item \textbf{העצמי הממשי (\textenglish{Actual Self})}: הייצוג של התכונות, היכולות, המראה והמעמד שהפרט מאמין שהוא מחזיק בהם בפועל.
    \item \textbf{העצמי האידיאלי (\textenglish{Ideal Self})}: הייצוג של התכונות וההישגים שהפרט שואף ומקווה להחזיק בהם -- מגלם משאלות לב ושאיפות עליונות.
    \item \textbf{העצמי הנדרש (\textenglish{Ought Self})}: הייצוג של החובות והאחריות שהפרט מאמין שעליו לקיים על פי נורמות חברתיות ומוסריות.
\end{enumerate}

פער חריף בין \textbf{העצמי הממשי} לבין \textbf{העצמי האידיאלי} גורר רגשות של דכדוך, עצב, אכזבה עמוקה ותחושת חוסר ערך (\textenglish{Dejection-Related Emotions}). במרחב הדיגיטלי, חשיפה רציפה למשפיענים המציגים עושר ויופי מושלם מגדירה מחדש את הסטנדרטים של "העצמי האידיאלי", והמרחק מ"העצמי הממשי" מתרחב משמעותית.

במקביל, תיאוריית ההחפצה העצמית (\textenglish{Self-Objectification Theory}) שהורחבה למדיה חברתית על ידי פרלוף (\textenglish{Perloff})~\cite{perloff2014social}, מסבירה כיצד חשיפה ממושכת לפלטפורמות מבוססות תמונה מעודדת משתמשים -- ובפרט נערות ונשים צעירות -- להפנים נקודת מבט של צופה חיצוני על גופן. תהליך זה של ניטור גופני מתמיד (\textenglish{Body Surveillance}) מרחיב את פער העצמי, גוזל משאבים קוגניטיביים יקרי ערך ומוביל לעלייה במדדי בושה וחרדה חברתית.

\section{חשיפה דיגיטלית, ערך עצמי ומצוקה רגשית: ראיות אמפיריות מרכזיות}
\label{sec:lit_empirical_evidence}
ההשפעות הפסיכולוגיות של מדיה חברתית נחקרו רבות בשני העשורים האחרונים. חלק זה מציג סקירה ביקורתית של מחקרי מפתח אמפיריים:
\begin{itemize}
    \item \textbf{מגמות אפידמיולוגיות ובריאות הנפש}: טוונגי ועמיתיה~\cite{twenge2018increases} הדגימו במחקר אורכי רחב היקף בקרב מעל 500,000 מתבגרים זינוק תלול ומובהק בדיכאון, מחשבות אובדניות ומצוקה רגשית מאז שנת 2010, במתאם הדוק עם חדירת הסמארטפונים והרשתות.
    \item \textbf{השלכות הגלילה הפסיבית לעומת אינטראקציה פעילה}: ורדואין וחוב'~\cite{verduyn2017social} הצביעו במטא-אנליזה על כך שצריכה פסיבית של תוכן היא הגורם המרכזי המתווך ירידה ברווחה דרך השוואה כלפי מעלה וקנאה, בעוד שאינטראקציה ישירה ואקטיבית עשויה לחזק הון חברתי.
    \item \textbf{בידוד חברתי נתפס}: פרימאק ועמיתיו~\cite{primack2017social} מצאו כי שימוש תדיר במגוון פלטפורמות מנבא פי 3 תחושת בידוד חברתי נתפס (\textenglish{PSI}).
    \item \textbf{אשליית האושר של האחר}: הטיות זמינות גורמות למשתמש להפריז בהערכת אושרם של אחרים ולהמעיט בהערכת מצוקותיהם, כפי שהדגימו ג'ורדן ועמיתיו~\cite{jordan2011misery}.
\end{itemize}

\section{דימוי גוף ושיווק משפיענים: מנגנוני שכנוע והחפצה}
\label{sec:lit_body_image_influencers}
אינסטגרם וטיקטוק הפכו לזירות מרכזיות של שיווק משפיענים. פרדולי וחוב'~\cite{fardouly2018instagram} הוכיחו כי חשיפה לתמונות משפיענים מגבירה השוואת מראה וחוסר שביעות רצון מגוף האישה. מחקריהם של לו ויוּאן (\textenglish{Lou \& Yuan})~\cite{lou2019influencer} ושל דה-וירמן וחוב' מנתחים את כוח השכנוע של משפיענים: (1)~\textbf{אינטראקציה פארא-סוציאלית} (\textenglish{Parasocial Interaction}) היוצרת תחושת אינטימיות כוזבת ומחלישה ספקנות צרכנית; (2)~\textbf{טשטוש תוכן ממומן} המשתלב בשגרת חיים לכאורה "אותנטית"; ו-(3)~\textbf{הוכחה חברתית} (\textenglish{Likes}, עוקבים) המקנה לגיטימציה מיידית לאידיאלים המוצגים.

\section{מודל מושגי רב-משתני של השוואה חברתית ברשתות}
\label{sec:lit_mediators_moderators}
איור~\ref{fig:mediator_moderator_model} פורס את המודל המושגי הרב-משתני שגובש ב-\projectname{}. הקשר בין חשיפה לפגיעה ברווחה מתווך על ידי ארבעה נתיבים קוגניטיביים: השוואה כלפי מעלה, פערי עצמי, אמון פארא-סוציאלי וחרדת החמצה (\textenglish{FOMO}). תיווך זה ממותן על ידי דפוס הגלילה (פסיבי מול אקטיבי) ועוצמת הסיוע המטא-קוגניטיבי של המערכת.

\begin{figure}[H]
\centering
\begin{LTR}
\begin{tikzpicture}[
    scale=0.62, every node/.style={transform shape},
    ivnode/.style={rectangle, draw=DarkPurple, fill=LavenderCard, rounded corners=3mm, line width=1.2pt, align=center, inner sep=6pt, text width=3.2cm, font=\small\sffamily},
    dvnode/.style={rectangle, draw=AlertRose, fill=AlertRose!15, rounded corners=3mm, line width=1.2pt, align=center, inner sep=6pt, text width=3.2cm, font=\small\sffamily},
    mednode/.style={rectangle, draw=RoyalPurple, fill=white, rounded corners=2.5mm, line width=1.0pt, align=center, inner sep=5pt, text width=3.4cm, font=\footnotesize\sffamily},
    modnode/.style={rectangle, draw=PurpleAccent, fill=white, rounded corners=2.5mm, line width=1.1pt, dashed, align=center, inner sep=5pt, text width=4.0cm, font=\footnotesize\sffamily},
    arrow/.style={-{Latex[length=2.5mm]}, line width=1.0pt, draw=DarkPurple},
    modarrow/.style={-{Latex[length=2.5mm]}, line width=1.0pt, draw=PurpleAccent, dashed}
]

\node[ivnode] (iv) {\textbf{Independent Variable}\\\scriptsize \textbf{Digital Content Exposure}\\\scriptsize High-Idealization Lifestyle,\\Influencer Sponsored Content};
\node[dvnode, right=9.2cm of iv] (dv) {\textbf{Dependent Variable}\\\scriptsize \textbf{Psychological Well-Being}\\\scriptsize State Self-Esteem ($S_t$),\\Affective Balance ($\mathbf{A}_t$),\\Body Image Dissatisfaction};

\node[mednode, right=2.6cm of iv, yshift=2.2cm] (med1) {\textbf{\textenglish{M1: Upward Comparison}}\\\scriptsize Contrast Effect Activation\\\scriptsize \textenglish{(Festinger 1954; Appel 2020)}};
\node[mednode, right=2.6cm of iv, yshift=0.7cm] (med2) {\textbf{\textenglish{M2: Self-Discrepancy}}\\\scriptsize Actual vs. Ideal Gap Expansion\\\scriptsize \textenglish{(Higgins 1987)}};
\node[mednode, right=2.6cm of iv, yshift=-0.7cm] (med3) {\textbf{\textenglish{M3: Parasocial Trust}}\\\scriptsize Disguised Commercial Persuasion\\\scriptsize \textenglish{(Lou \& Yuan 2019)}};
\node[mednode, right=2.6cm of iv, yshift=-2.2cm] (med4) {\textbf{\textenglish{M4: Fear of Missing Out}}\\\scriptsize FOMO \& Social Exclusion\\\scriptsize \textenglish{(Primack et al. 2017)}};

\begin{scope}[on background layer]
    \node[draw=LavenderBorder, fill=LavenderLight, rounded corners=3mm, fit=(med1) (med4), inner sep=6pt, label={[DarkPurple, font=\small\bfseries\sffamily]above:\textenglish{Mediating Cognitive Pathways}}] (med_group) {};
\end{scope}

\draw[arrow] (iv.north east) -- (med1.west);
\draw[arrow] (iv.east) -- (med2.west);
\draw[arrow] (iv.east) -- (med3.west);
\draw[arrow] (iv.south east) -- (med4.west);

\draw[arrow] (med1.east) -- (dv.north west);
\draw[arrow] (med2.east) -- (dv.west);
\draw[arrow] (med3.east) -- (dv.west);
\draw[arrow] (med4.east) -- (dv.south west);

\draw[-{Latex[length=2.2mm]}, dotted, line width=1.0pt, draw=SlateGray] (iv.south) -- ++(0,-1.3cm) -| node[below, pos=0.25, font=\scriptsize\sffamily, text=SlateGray] {\textenglish{Direct Effect (Unmediated Residual)}} (dv.south);

\node[modnode, above=1.1cm of med_group] (mod_top) {\textbf{\textenglish{Moderator 1: Interaction Mode}}\\\scriptsize Passive Scrolling vs. Active Direct Chat\\\scriptsize \textenglish{(Verduyn et al. 2017)}};
\node[modnode, below=1.1cm of med_group] (mod_bottom) {\textbf{\textenglish{Moderator 2: Metacognitive Insight}}\\\scriptsize \projectname{} Cognitive Mirror Scaffolding\\\scriptsize \textenglish{(Buffers \& Neutralizes Impact)}};

\draw[modarrow] (mod_top.south) -- (med_group.north);
\draw[modarrow] (mod_bottom.north) -- (med_group.south);
\end{tikzpicture}
\end{LTR}
\caption{מודל מושגי רב-משתני: שרשרת התיווך הקוגניטיבית והמשתנים הממתנים בין חשיפה דיגיטלית לרווחה נפשית}
\label{fig:mediator_moderator_model}
\end{figure}

\section{כלי מדידה פסיכומטריים והמרתם למשתני מצב חישוביים}
\label{sec:lit_psychometric_instruments}
כדי שמערכת מרובת סוכנים תוכל למדל, לנתח ולשקף מצבים רגשיים, עליה להשתית את משתני המצב הפנימיים על מכשירי מדידה פסיכומטריים בעלי מהימנות ותוקף מדעי מוכחים:
\begin{enumerate}
    \item \textbf{סולם הערך העצמי של רוזנברג (\textenglish{Rosenberg Self-Esteem Scale - RSES})}~\cite{rosenberg1965society}: המדד המתוקף ביותר להערכת ערך עצמי כללי ורגעי (10 פריטים, סולם ליקרט 1--4). במערכת \projectname{}, שינויים רגעיים בערך העצמי ממוסגרים סביב ציון מנורמל $S_t \in [0, 1]$.
    \item \textbf{לוח הרגשות החיוביים והשליליים (\textenglish{PANAS})}~\cite{watson1988development}: שאלון הכולל 20 תיאורי רגש (10 חיוביים ו-10 שליליים) המודדים מצבי רגש זמניים בעלי עוררות גבוהה $\mathbf{A}_t = (\mathrm{PA}_t, \mathrm{NA}_t)$.
    \item \textbf{סולמות הרווחה הפסיכולוגית של ריף (\textenglish{Ryff's Scales of Psychological Well-Being})}~\cite{ryff1989happiness}: מודל רב-ממדי הבוחן שישה ממדים מרכזיים של בריאות נפשית חיובית: אוטונומיה, שליטה בסביבה, צמיחה אישית, יחסים חיוביים, מטרה בחיים וקבלה עצמית ($\mathbf{R}_t \in [0, 1]^6$).
    \item \textbf{מדדי דימוי גוף ורזון (\textenglish{EDI / SATAQ})}~\cite{perloff2014social}: הפנמת מודל אידיאל יופי וחוסר שביעות רצון גופני המשוקללים במדד החשיפה $\exposureIndex{} \in [0, 1]$.
\end{enumerate}

טבלה~\ref{tab:psychometric_mapping} מציגה ניתוח השוואתי של כלי מדידה אלו ואת אופן המרתם למשתני מצב וקטוריים במערכת \projectname{}.

\begin{table}[H]
\centering
\caption{מיפוי כלים פסיכומטריים מתוקפים למשתני מצב חישוביים במערכת \projectname{}}
\label{tab:psychometric_mapping}
\footnotesize
\renewcommand{\arraystretch}{0.90}
\begin{tabularx}{\textwidth}{>{\raggedleft\arraybackslash}p{2.6cm} >{\raggedleft\arraybackslash}p{2.5cm} >{\raggedleft\arraybackslash}p{3.2cm} X}
\toprule
\rowcolor{LavenderCard}
\textbf{\color{DarkPurple}כלי פסיכומטרי} & \textbf{\color{DarkPurple}ממדים נמדדים} & \textbf{\color{DarkPurple}משתנה חישובי במערכת} & \textbf{\color{DarkPurple}משמעות פונקציונלית בארכיטקטורה} \\
\midrule
\textbf{\textenglish{RSES}} \newline \textenglish{Rosenberg (1965)}~\cite{rosenberg1965society} & 
הערכה עצמית גלובלית ורגעי (10 פריטים) & 
$S_t \in [0, 1]$ \newline רכיב הערך העצמי ב-$\wellbeingVector{}$ & 
כיול ירידה בערך העצמי עקב אירועי השוואה כלפי מעלה בעלי $\comparisonScore{}$ גבוה. \\
\midrule
\textbf{\textenglish{PANAS}} \newline \textenglish{Watson et al. (1988)}~\cite{watson1988development} & 
מאזן אפקטיבי: רגש חיובי (\textenglish{PA}) מול שלילי (\textenglish{NA}) & 
$\mathbf{A}_t = (\mathrm{PA}_t, \mathrm{NA}_t)$ \newline רכיב אפקטיבי דו-קוטבי & 
ניטור תנודות במצב הרוח ועלייה בעוררות שלילית לאחר גלילה פסיבית מתמשכת. \\
\midrule
\textbf{\textenglish{Ryff's PWB}} \newline \textenglish{Ryff (1989)}~\cite{ryff1989happiness} & 
רווחה ארוכת-טווח: אוטונומיה, שליטה, קבלה עצמית & 
$\mathbf{R}_t \in [0, 1]^6$ \newline וקטור רווחה מצטבר & 
שיקוף תובנות מאקרו על ידי \agentReflector{} בסיכומי שימוש שבועיים. \\
\midrule
\textbf{\textenglish{EDI / SATAQ}} \newline דימוי גוף ורזון~\cite{perloff2014social} & 
הפנמת אידיאל יופי, חוסר שביעות רצון גופני & 
$\exposureIndex{} \in [0, 1]$ \newline מדד חשיפה למשפיענים & 
סיווג רמת הסיכון של תכני יופי ולייפסטייל על ידי \agentAnalyst{}. \\
\bottomrule
\end{tabularx}
\end{table}

\section{סוכנים גנרטיביים כמודל לסימולציה אנושית (Park et al., 2023)}
\label{sec:lit_flagship_park2023}
כדי לגשר בין תיאוריות פסיכולוגיות למידול חישובי אמין של התנהגות משתמשים ברשתות, פרויקט \projectname{} מאמץ ומפתח את הפרדיגמה החישובית שהוצגה במאמר הדגל של פארק ועמיתיו (\textenglish{Park et al., 2023}) מאוניברסיטת סטנפורד ו-\textenglish{Google Research}~\cite{park2023generative}. במחקרם, יושמה סביבת ארגז חול דיגיטלית (\textenglish{Smallville}) המאוכלסת על ידי 25 סוכני פרסונה גנרטיביים אוטונומיים. שופטים אנושיים דירגו את רמת האמינות (\textenglish{Believability}) של התנהגות הסוכנים כעולה באופן מובהק על מודלי בסיס ($p < .001$), בהקשרים של הבנה חברתית ותגובות רגשיות עקביות.

ארכיטקטורת הסוכן ב-\textenglish{Smallville} מושתתת על שלושה רכיבים חישוביים מרכזיים: \textbf{זיכרון שלם} (\textenglish{Memory Stream}), \textbf{רפלקציה} (\textenglish{Reflection} -- סינתזה היררכית של אירועים לתובנות מופשטות), ו\textbf{תכנון} (\textenglish{Planning}). \projectname{} משאיל את מנגנוני הסינתזה והרפלקציה של פארק וחוב' כדי למדל את תהליכי שחיקת הרווחה בעת גלילה פסיבית, ומיישם מנגנון שאיפה-רפלקציה מקומי בסוכן הרפלקציה (\agentReflector{}).

\section{מטריצה סינתטית של הספרות האמפירית המרכזית}
\label{sec:lit_summary_table}
טבלה~\ref{tab:empirical_literature_matrix} מסכמת את גופי הידע המרכזיים שנסקרו בפרק זה ואת השלכותיהם הישירות על התכן ההנדסי של \projectname{}.

\begin{table}[H]
\centering
\caption{מטריצה סינתטית של הספרות האמפירית המרכזית והשלכותיה לתכן \projectname{}}
\label{tab:empirical_literature_matrix}
\footnotesize
\renewcommand{\arraystretch}{0.88}
\begin{tabularx}{\textwidth}{>{\bfseries\raggedleft\arraybackslash}p{2.6cm} >{\raggedleft\arraybackslash}p{2.3cm} >{\raggedleft\arraybackslash}p{4.2cm} X}
\toprule
\rowcolor{LavenderCard}
\textbf{\color{DarkPurple}מחקר ומקור} & \textbf{\color{DarkPurple}מתודולוגיה ומדגם} & \textbf{\color{DarkPurple}ממצאים אמפיריים מרכזיים} & \textbf{\color{DarkPurple}השלכה הנדסית ל-\projectname{}} \\
\midrule
\textbf{\textenglish{Festinger (1954)}}~\cite{festinger1954theory} & מודל תיאורטי קלאסי & דחף השוואה מולד; כיווניות מעלה/מטה; תנאי ניגוד והטמעה. & בסיס לפונקציית ציון בולטות ההשוואה $\comparisonScore{}$. \\
\midrule
\textbf{\textenglish{Higgins (1987)}}~\cite{higgins1987self} & ניסויי ותיאורטי & פערי עצמי מנבאים ישירות פגיעות לרגשות דכדוך ואובדן ערך. & הגדרת רכיבי וקטור המצב $\wellbeingVector{}$ ומדדי פער עצמי. \\
\midrule
\textbf{\textenglish{Twenge et al. (2018)}}~\cite{twenge2018increases} & אורכי ($N > 500{,}000$) & זינוק חריף בדיכאון ובמצוקת נוער מאז 2010 במתאם לחדירת סמארטפונים. & תיקוף הצורך בהתערבות ואפיון אוכלוסיית היעד. \\
\midrule
\textbf{\textenglish{Primack et al. (2017)}}~\cite{primack2017social} & חתך לאומי ($N = 1{,}787$) & שימוש בריבוי רשתות מנבא פי 3 בידוד חברתי נתפס (\textenglish{PSI}). & ניטור פיזור שימוש ורצפי אפליקציות על ידי \agentObserver{}. \\
\midrule
\textbf{\textenglish{Verduyn et al. (2017)}}~\cite{verduyn2017social} & סקירה ומטא-אנליזה & גלילה פסיבית מתווכת ירידה ברווחה דרך השוואה כלפי מעלה וקנאה. & הבחנה קריטית בין שימוש פסיבי לאקטיבי באלגוריתם הניתוח. \\
\midrule
\textbf{\textenglish{Fardouly et al. (2018)}}~\cite{fardouly2018instagram} & ניסויי-מתאמי ($N = 227$) & תמונות משפיענים מעוררות החפצה עצמית ואי-שביעות רצון גופנית. & שילוב סיווג תכני יופי ולייפסטייל במדד $\exposureIndex{}$. \\
\midrule
\textbf{\textenglish{Appel et al. (2015, 2016)}}~\cite{appel2015social, appel2016interplay} & שדה ומעבדה מבוקרים & אישוש אפקט הניגוד: חשיפה לסטנדרטים אידיאליים מנבאת קנאה ושחיקה. & כיול קצב השחיקה האקספוננציאלי של מדד הערך העצמי הרגעי. \\
\midrule
\textbf{\textenglish{Lou \& Yuan (2019)}}~\cite{lou2019influencer} & משוואות מבניות ($N = 538$) & אינטראקציה פארא-סוציאלית מגבירה אמון ומחלישה ביקורתיות לפרסום. & זיהוי טכניקות מיתוג סמוי ושיווק ממומן על ידי \agentAnalyst{}. \\
\midrule
\textbf{\textenglish{Park et al. (2023)}}~\cite{park2023generative} & סימולציית סוכני LLM & סוכנים גנרטיביים מבוססי זיכרון ורפלקציה מפגינים אמינות חברתית. & ארכיטקטורת מודל הסימולציה ומנגנון הרפלקציה המקומי. \\
\bottomrule
\end{tabularx}
\end{table}

\begin{insightbox}{תובנה מחקרית מרכזית מתוך סקירת הספרות}
הקונצנזוס המדעי ברור: הנזק הפסיכולוגי ברשתות חברתיות אינו פועל יוצא של משך הזמן הכולל, אלא תוצאה ישירה של \textbf{איכות האינטראקציה, אופי התוכן המוצג (תכני לייפסטייל ומשפיענים אידיאליים) ועירור תהליכי השוואה כלפי מעלה בעת גלילה פסיבית}. ממצא זה מציב את הדרישה המרכזית לפלטפורמת \projectname{}: מעבר מניהול זמן גס לניתוח סמנטי-איכותני של גירויי השוואה באמצעות סוכנים אוטונומיים.
\end{insightbox}
